% ============================================================
%  CHAPITRE 2 — Section 5 : Intégration réactive du vieillissement
%                            des composants
% ============================================================

\section{Intégration réactive du vieillissement des composants}

Les résultats préliminaires présentés dans la section précédente ont permis d'identifier la stratégie RB2 comme le meilleur point de départ pour une extension adaptative. L'idée directrice de cette section est d'aller au-delà d'une stratégie statique, dont les paramètres sont figés une fois pour toutes, en introduisant une couche d'adaptation qui ajuste dynamiquement le comportement de l'EMS en fonction de l'état de santé courant des composants. Cette approche, qualifiée de \textit{réactive au vieillissement} (\textit{health-conscious}), vise à tirer parti de l'information disponible sur le SoH pour modifier les décisions de contrôle au fil du temps, sans recourir à des méthodes d'optimisation complexes ni à des prévisions de l'état futur du système.

% ------------------------------------------------------------
\subsection{Principe de l'adaptation par le SoH}
% ------------------------------------------------------------

Le vieillissement des composants électrochimiques se traduit concrètement par une dégradation de leurs performances~: la capacité utile de la batterie diminue, la puissance délivrable par la pile à combustible à courant constant se réduit, et la puissance consommée par l'électrolyseur à courant constant augmente. Ces évolutions sont capturées par les indicateurs de SoH définis à la section~2. Une conséquence directe de ce vieillissement est que les plages de fonctionnement nominales — définies lors de la mise en service du système — deviennent progressivement inadaptées~: continuer à solliciter un composant dégradé aux mêmes niveaux de puissance que lorsqu'il était neuf revient à l'opérer dans des conditions de plus en plus contraignantes, ce qui accélère davantage sa dégradation.

L'idée centrale de l'adaptation par le SoH est donc de faire évoluer les paramètres de contrôle en cohérence avec l'état réel des composants, en réduisant progressivement les niveaux de puissance demandés à mesure que ceux-ci vieillissent. Cette réduction est naturelle~: un composant qui a perdu une fraction de ses performances ne peut physiquement plus délivrer la même puissance qu'à l'état neuf, et il est souhaitable que l'EMS en tienne compte explicitement plutôt que de tenter de compenser cette perte par une sollicitation accrue. En pratique, cela revient à introduire une boucle de retour lente sur le SoH, qui agit sur les paramètres de la stratégie à une échelle de temps compatible avec celle du vieillissement — typiquement à la semaine ou au mois — sans perturber le contrôle temps réel.

Cette approche se distingue des méthodes prédictives ou d'optimisation en ligne~: elle ne nécessite ni modèle de prévision du SoH futur, ni résolution d'un problème d'optimisation à chaque pas de temps. Elle repose uniquement sur la valeur courante du SoH estimé, ce qui la rend simple, robuste et directement implémentable dans un système embarqué.

% ------------------------------------------------------------
\subsection{Application à la stratégie RB2}
% ------------------------------------------------------------

La stratégie RB2, rappelons-le, repose sur l'allocation d'une puissance d'offset fixe aux composants hydrogène~: $P_{\text{FC-offset}}$ pour la pile à combustible lorsque la demande excède la production, et $P_{\text{ELY-offset}}$ pour l'électrolyseur dans le cas contraire. La batterie assure ensuite le complément. Ces offsets sont calibrés une fois pour toutes lors de la phase de conception, en supposant que les composants sont à l'état neuf.

L'adaptation par le SoH consiste à moduler ces offsets en fonction de l'état de santé estimé des composants hydrogène. La logique est la suivante~: lorsqu'un composant se dégrade, sa puissance maximale délivrable diminue, et il est cohérent de réduire proportionnellement la puissance qui lui est demandée. Cela maintient les composants dans des plages de fonctionnement relatives similaires tout au long de leur vie, évitant ainsi de les contraindre à opérer dans des régimes de plus en plus extrêmes à mesure qu'ils vieillissent.

% - - - - - - - - - - - - - - - - - - - - - - - - - - - - - -
\subsubsection{Adaptation de la puissance de la PAC}
% - - - - - - - - - - - - - - - - - - - - - - - - - - - - - -

Lorsque la demande de charge excède la production PV ($\hat{P}_{\text{Load}} > \hat{P}_{\text{PV}}$), la puissance de consigne allouée à la pile à combustible est modulée par son SoH courant~:

\begin{equation}
    \tilde{P}_{\text{FC}} = P_{\text{FC-offset}} \cdot \text{SoH}_{\text{FC}}
    \label{eq:adaptation_fc}
\end{equation}

Ainsi, lorsque la pile à combustible est neuve ($\text{SoH}_{\text{FC}} = 1$), la puissance demandée est identique à celle de la stratégie RB2 de base. À mesure que la pile vieillit et que son SoH diminue, la puissance qui lui est allouée se réduit proportionnellement. Cette réduction traduit physiquement le fait que la pile est de moins en moins capable de délivrer des niveaux de puissance élevés sans subir de contraintes supplémentaires, et qu'il est préférable de l'opérer dans une plage plus confortable. La batterie, dont la dynamique est plus favorable, absorbe alors le déficit de puissance correspondant.

% - - - - - - - - - - - - - - - - - - - - - - - - - - - - - -
\subsubsection{Adaptation de la puissance de l'électrolyseur}
% - - - - - - - - - - - - - - - - - - - - - - - - - - - - - -

De manière symétrique, lorsque la production PV excède la demande de charge ($\hat{P}_{\text{Load}} < \hat{P}_{\text{PV}}$), la puissance de consigne allouée à l'électrolyseur est modulée par son SoH~:

\begin{equation}
    \tilde{P}_{\text{ELY}} = P_{\text{ELY-offset}} \cdot \text{SoH}_{\text{ELY}}
    \label{eq:adaptation_ely}
\end{equation}

Le raisonnement est identique~: un électrolyseur vieillissant consomme davantage de puissance électrique pour produire la même quantité d'hydrogène à courant constant. Lui demander des niveaux de puissance d'absorption élevés revient donc à le contraindre à des régimes de plus en plus défavorables. En réduisant progressivement la puissance qui lui est allouée en proportion de son SoH, l'EMS adaptative maintient l'électrolyseur dans des conditions d'utilisation plus clémentes, ralentissant ainsi son vieillissement futur. L'excès de production PV non absorbé par l'électrolyseur est alors redirigé vers la batterie, ou écrêté si celle-ci est également pleine.

Ces deux équations peuvent être regroupées de façon compacte~:

\begin{equation}
\label{eq:fc-ely-compact}
\begin{cases}
\tilde{P}_{\text{FC}}  = P_{\text{FC-offset}}  \cdot \text{SoH}_{\text{FC}}  & \text{si } \hat{P}_{\text{Load}} > \hat{P}_{\text{PV}} \\
\tilde{P}_{\text{ELY}} = P_{\text{ELY-offset}} \cdot \text{SoH}_{\text{ELY}} & \text{si } \hat{P}_{\text{Load}} < \hat{P}_{\text{PV}}
\end{cases}
\end{equation}

La stratégie ainsi obtenue est notée RB2(SoH) dans la suite du document.

% ------------------------------------------------------------
\subsection{Résultats \& analyse}
% ------------------------------------------------------------

La stratégie RB2(SoH) a été simulée dans les mêmes conditions que les EMS présentées à la section précédente, sur un horizon de quinze ans. Son positionnement dans le plan de Pareto est représenté à la Figure~\ref{fig:pareto}, et l'évolution temporelle des principales grandeurs du système est illustrée à la Figure~\ref{fig:curves}. Les quantités de vieillissement — évolution du SoH et répartition des types de dégradation pour les composants hydrogène — sont présentées à la Figure~\ref{fig:aging}.

\begin{figure}[h!]
    \centering
    \includegraphics[width=\linewidth]{Chapitre 2/SOH_feedback/Figures/everything_combined_v2_2.pdf}
    \caption{Profils de puissance sur quinze ans (gauche) et zoom sur une semaine (droite), EMS RB2(SoH).}
    \label{fig:curves}
\end{figure}
\FloatBarrier

\begin{figure}[h!]
    \centering
    \includegraphics[width=\linewidth]{Chapitre 2/SOH_feedback/Figures/all_aging_2.pdf}
    \caption{Évolution des grandeurs de vieillissement~: SoH et types de dégradation des composants hydrogène.}
    \label{fig:aging}
\end{figure}
\FloatBarrier

Il ressort clairement de la Figure~\ref{fig:pareto} que l'introduction du retour SoH permet de déplacer significativement le point de fonctionnement de RB2 dans le plan de Pareto, dans une direction favorable~: l'indicateur de dégradation est réduit d'environ $34{,}0\,\%$, au prix d'une légère augmentation de la LPSP d'environ $12{,}8\,\%$. Ce résultat confirme que l'adaptation réactive au vieillissement permet effectivement d'améliorer le compromis entre continuité de service et durabilité des composants, sans recourir à des méthodes d'optimisation complexes.

L'examen des profils temporels (Figure~\ref{fig:curves}) permet de comprendre le mécanisme sous-jacent. Sur l'horizon de quinze ans, on observe que les puissances allouées à la pile à combustible et à l'électrolyseur diminuent progressivement au fil du temps, en cohérence avec la réduction de leur SoH respectif. À titre d'illustration, la puissance nominale de la pile à combustible passe d'environ $0{,}95$\,kW en début de simulation à $0{,}86$\,kW en fin de simulation, et celle de l'électrolyseur de $-9{,}00$\,kW à $-8{,}10$\,kW. La batterie compense naturellement cette réduction, ce qui se traduit par une légère augmentation de sa sollicitation et, in fine, par la hausse modérée de la LPSP observée.

L'analyse des types de dégradation (Figure~\ref{fig:aging}) révèle par ailleurs que les cycles de démarrage-arrêt constituent de loin le facteur de dégradation le plus significatif pour les deux composants hydrogène. Pour la pile à combustible, cette contribution représente environ $94{,}5\,\%$ de la dégradation totale, les transitoires de puissance comptant pour les $5{,}5\,\%$ restants. Pour l'électrolyseur, la part des démarrages-arrêts est de $56{,}6\,\%$, les $43{,}4\,\%$ restants étant répartis entre le fonctionnement à forte puissance ($33{,}2\,\%$) et les transitoires ($10{,}2\,\%$). Des tests ont été réalisés pour évaluer si le maintien en mode ralenti (\textit{idling}) plutôt que l'arrêt complet lors des périodes de faible demande permettrait de réduire ces dommages. Pour la PEMFC, les résultats montrent que le maintien en mode ralenti tend à détériorer les performances plutôt qu'à les améliorer~; pour le PEMWE, les résultats ne sont pas meilleurs qu'en cas d'arrêt complet. Ces résultats suggèrent que la réduction des cycles de démarrage-arrêt devrait être une priorité dans les futures extensions de ce travail.

% ------------------------------------------------------------
\subsection{Discussion : compromis dégradation / LPSP}
% ------------------------------------------------------------

Les résultats obtenus mettent en évidence un compromis fondamental inhérent à toute EMS réactive au vieillissement~: la préservation des composants se fait nécessairement au détriment d'une certaine dégradation de la qualité de service. Dans le cas de RB2(SoH), ce compromis se matérialise par une réduction de $34{,}0\,\%$ des coûts de remplacement sur quinze ans, contre une augmentation de $12{,}8\,\%$ de la LPSP par rapport à RB2. Ce résultat n'est pas surprenant~: en réduisant la puissance disponible auprès des composants hydrogène, on diminue la capacité totale du système à répondre à la demande dans les situations où la batterie seule ne suffit pas.

Il est important de souligner que ce compromis n'est pas figé. La formulation adoptée ici, dans laquelle les offsets sont multipliés directement par le SoH, constitue le cas particulier d'une loi d'adaptation linéaire. D'autres formes de modulation sont envisageables — par exemple des lois affines, exponentielles ou par paliers — qui permettraient d'ajuster la sévérité de la réduction en fonction du niveau de dégradation atteint, ou de définir un seuil de SoH en dessous duquel l'adaptation s'active. Le choix d'une telle loi constitue un levier de conception supplémentaire pour l'opérateur, qui peut ainsi cibler un point de fonctionnement souhaité dans le plan de Pareto.

Par ailleurs, les résultats présentés dans cette section sont intrinsèquement liés aux profils de puissance utilisés et au dimensionnement du système. Une analyse de sensibilité à ces deux paramètres permettrait d'évaluer dans quelle mesure les conclusions obtenues sont généralisables à d'autres contextes. De même, la prise en compte d'échelles de temps intermédiaires — robustesse face aux défaillances temporaires de composants, par exemple — constitue une perspective d'extension naturelle du cadre proposé, dont les bases ont été jetées dans la section~3. Ces limitations et perspectives seront reprises dans le cadre du bilan général du chapitre.

\section{Bouclage avec le dimensionnement du micro-réseau}
\lipsum[1]